\documentclass[11pt]{article}
\usepackage{icmlko}
\begin{document}

\title{표본을 두 배 넘게 늘렸더니 답이 더 흐려졌다 --- 이웃 효과는 미리 잴 수 없다}
\author{월드모델 실험실}
\date{노트 599 · 600 · 2026년 8월 5일}
\maketitle

\begin{abstract}
논문 $457$ 이 스스로 적었다 --- ``\textbf{이웃 효과를 여섯 번 만났는데
매번 사후에만 안다. 그것을 미리 재는 자가 없다.}'' 예측기를 만들려 했더니
\textbf{전역 변경과 단일 도메인 변경을 둘 다 잰 짝이 셋뿐}이었다(노트
$599$). 넷을 더 만들어 \textbf{일곱}으로 늘렸다.

\begin{center}
\begin{tabular}{lrr}
\hline
 & 짝 $3$(노트 $599$) & \textbf{짝 $7$(노트 $600$)} \\
\hline
이웃 차 범위 & $0.036 \sim 0.128$ & $\mathbf{-0.028 \sim 0.128}$ \\
부호 뒤집힘 & $1/3$ & $\mathbf{4/7}$ \\
덮음과의 순위상관 & --- & $\mathbf{+0.107}$ \\
\hline
\end{tabular}
\end{center}

\textbf{표본을 두 배 넘게 늘렸는데 답이 더 흐려졌다.} 부호가 뒤집히는
짝이 셋 중 하나에서 \textbf{일곱 중 넷}이 됐고, 덮음은 이웃 효과를
\textbf{전혀} 예측하지 못한다($|$차$|$ 기준 $+0.036$).
\end{abstract}

\begin{figure}[t]
\centering
\includegraphics[width=\textwidth]{figs/f1.pdf}
\caption{\textbf{가} --- 전역에서 단일로 가면 절반 넘게 부호가 바뀐다.
\textbf{나} --- 덮음은 예측하지 못한다.
\textbf{다} --- 짝을 늘렸더니 나빠졌다.}
\end{figure}

\section{일곱 짝}

\begin{center}
\begin{tabular}{lrrrr}
\hline
도메인 & 덮음 & 전역 차 & 단일 차 & 이웃 차 \\
\hline
도서 & $69\%$ & $+0.0459$ & $\mathbf{-0.0820}$ & $\mathbf{+0.1279}$ \\
팝업 & $\mathbf{0\%}$ & $+0.0374$ & $+0.0018$ & $+0.0356$ \\
게임 & $67\%$ & $-0.0137$ & $\mathbf{+0.0138}$ & $-0.0275$ \\
웹툰 & $76\%$ & $-0.0213$ & $-0.0285$ & $+0.0072$ \\
모바일 & $80\%$ & $+0.0005$ & $\mathbf{-0.0062}$ & $+0.0067$ \\
세계애니 & $56\%$ & $+0.0039$ & $\mathbf{-0.0027}$ & $+0.0066$ \\
애니 & $64\%$ & $-0.0004$ & $+0.0000$ & $-0.0004$ \\
\hline
\end{tabular}
\end{center}

\textbf{덮음 $0\%$ 인 팝업의 이웃 차가 $+0.0356$ 이고, 덮음 $80\%$ 인
모바일이 $+0.0067$ 이다.} 덮음이 높으면 직접 몫이 커서 이웃 몫이 작을
것 같지만 \emph{그렇지 않다} --- 도서는 덮음 $69\%$ 인데 이웃 차가
$+0.1279$ 로 제일 크다.

\section{내 조항에 내가 걸렸다}

처음 낸 표는 이웃 몫을 \textbf{비}로 적었고, 모바일이 $\mathbf{1228\%}$
였다. 전역 차가 $+0.0005$ --- \textbf{분모가 $0$ 근처}다.

\begin{quote}
\textbf{노트 $584$ 가 세운 조항이 정확히 이것이다} --- ``비가 $3$ 을 넘으면
기준선이 $0$ 근처다. \textbf{비 대신 차를 적는다}''(\texttt{sideaudit.
baseline\_ratio}). 내가 만든 경고가 겨눈 실수를 \emph{다른 표}에서
그대로 했다.
\end{quote}

차로 다시 적으니 모바일은 $+0.0067$ 로 웹툰($+0.0072$) $\cdot$
세계애니($+0.0066$)와 \textbf{같은 급}이다. 그리고 \textbf{결론이 더
강해진다} --- 덮음과의 순위상관이 비 기준 $+0.429$ 에서 차 기준
$\mathbf{+0.107}$ 로 내려간다.

\section{그래서 무엇이 남나}

\begin{itemize}
\item \textbf{``전역 표로 후보를 세우지 않는다''}(노트 $551 \cdot 560$)가
$n{=}7$ 로 확정됐다. 부호가 \textbf{절반 넘게} 뒤집힌다.
\item \textbf{크기도 부호도 예측할 수 없다.} 유일하게 확실한 것은
\emph{덮음 $0\%$ 면 차가 전부 이웃 효과}인데, 그것은 \textbf{정의상
참}이지 예측이 아니다.
\item \texttt{sideaudit.domain\_deltas} 가 지금 형태로 옳다 --- 덮음 칸을
찍고, $0\%$ 인 줄을 표시하고, \textbf{나머지는 ``못 잰다''고 적는다.}
\end{itemize}

\section*{조항}

\begin{quote}
\textbf{표본을 늘려도 답이 안 나오면 그것이 답이다.} 셋에서 일곱으로
늘렸는데 범위가 넓어지고 부호 뒤집힘이 $1/3$ 에서 $4/7$ 이 됐다.
\textbf{``더 재면 알 것''을 무한히 미루지 않으려면, 늘려 보고 나빠졌다는
것을 적어야 한다.}
\end{quote}

\section*{남은 것}

\begin{itemize}
\item \textbf{손잡이가 하나뿐}이다(\texttt{zero\_is\_data}). 다른 계열
(축 가림 $\cdot$ 달력)에서도 같은지 안 봤다 --- 다만 옛 짝 셋도 같은
손잡이였으므로 \textbf{일곱이 전부 한 계열}이다.
\item \textbf{애니는 표에서 빼야 했다} --- 전역 $-0.0004$ $\cdot$ 단일
$+0.0000$ 으로 둘 다 $0$ 이다. 이미 \texttt{DOMDROP} 으로 검색을 뗀
도메인이라 \emph{바꿀 것이 없었다}. 짝이 실은 여섯이다.
\item \textbf{원장의 옛 수치도 비로 적혀 있다} --- $95\% \cdot 279\% \cdot
199\%$. 차로는 $+0.0356 \cdot +0.1279 \cdot -0.0275$ 다.
\end{itemize}

\end{document}
